\documentclass[12pt]{article}
\usepackage{graphicx}
\usepackage{geometry}
\usepackage{float}
\geometry{margin=1in}

\begin{document}

%================ HEADER =================%

\begin{center}

\includegraphics[width=2.5cm]{iiitb_comet_logo.jpg}

\vspace{0.2cm}

\textbf{Name: Srinidhi A} \\
\textbf{ID: COMETFWC053}
\end{center}

\vspace{0.3cm}

\begin{center}
\huge \textbf{GATE Question No. 41}
\end{center}

\vspace{0.3cm}

%================ QUESTION =================%

\textbf{Question:}

For any set of inputs, $A$ and $B$, the following circuits give the same output, $Q$, except one. Identify which one is different.

\vspace{0.2cm}

% Image forced to stay here
\begin{figure}[H]
\centering
\includegraphics[width=0.6\textwidth]{question.jpg}
\end{figure}

\begin{center}
(GATE PH 2010)
\end{center}

\vspace{0.4cm}

%================ SOLUTION =================%

\huge \textbf{Solution}

\normalsize

We simplify each circuit using Boolean algebra.

\vspace{0.4cm}

\textbf{Step 1: Circuit (a)}

From the diagram, output is AND of $(A + B)$ and $B'$.

$Q_a = (A + B)\cdot B'$

Using distributive law:

$Q_a = A B' + B B'$

Since $B B' = 0$,

$Q_a = A B'$

\vspace{0.4cm}

\textbf{Step 2: Circuit (b)}

$Q_b = A B'$

\vspace{0.4cm}

\textbf{Step 3: Circuit (c)}

After simplifying internal gates,

$Q_c = A B'$

\vspace{0.4cm}

\textbf{Step 4: Circuit (d)}

$Q_d = A' + B$

\vspace{0.6cm}

\textbf{Truth Table:}

\begin{center}
\begin{tabular}{|c|c|c|c|}
\hline
A & B & $A B'$ & $A' + B$ \\
\hline
0 & 0 & 0 & 1 \\
0 & 1 & 0 & 1 \\
1 & 0 & 1 & 0 \\
1 & 1 & 0 & 1 \\
\hline
\end{tabular}
\end{center}

\vspace{0.6cm}

Circuits (a), (b), and (c) give same output.

Circuit (d) gives different output.

\vspace{0.8cm}

\begin{center}
\huge \boxed{\text{Correct Answer: Option (d)}}
\end{center}

\end{document}
